\textbf{Phần 1. Trắc nghiệm lựa chọn}
\setcounterpageref{ex}{0}

\begin{ex} %[2D2N14]
Kết quả của phép tính $9 + 6$ là
\choice
{\True $15$}
{$14$}
{$16$}
{$17$}
\loigiai{
$9 + 6 = 15$.\\
\textbf{Đáp án của câu hỏi là A}
}
\end{ex}

\begin{ex} %[2D2N14]
Kết quả của phép tính $15 - 7$ là
\choice
{$7$}
{$9$}
{$10$}
{\True $8$}
\loigiai{
$15 - 7 = 8$.\\
\textbf{Đáp án của câu hỏi là D}
}
\end{ex}

\begin{ex} %[2D2H14]
Kết quả của phép tính $8 + 7$ là
\choice
{\True $15$}
{$14$}
{$16$}
{$17$}
\loigiai{
$8 + 7 = 15$.\\
\textbf{Đáp án của câu hỏi là A}
}
\end{ex}

\begin{ex} %[2D2H14]
Kết quả của phép tính $17 - 9$ là
\choice
{$7$}
{$9$}
{\True $8$}
{$10$}
\loigiai{
$17 - 9 = 8$.\\
\textbf{Đáp án của câu hỏi là C}
}
\end{ex}

\begin{ex} %[2D2H14]
Điền dấu thích hợp: $13 - 5 \dots 14 - 6$.
\choice
{$>$}
{\True $=$}
{$<$}
{Không so sánh được}
\loigiai{
$13 - 5 = 8$ và $14 - 6 = 8$. Vậy $13 - 5 = 14 - 6$.\\
\textbf{Đáp án của câu hỏi là B}
}
\end{ex}

\begin{ex} %[2D2H14]
Tính giá trị biểu thức: $16 - 8 + 5$.
\choice
{\True $13$}
{$12$}
{$14$}
{$15$}
\loigiai{
$8 + 5 = 13$.\\
\textbf{Đáp án của câu hỏi là A}
}
\end{ex}

\begin{ex} %[2D2V14]
Một cửa hàng buổi sáng bán được 8 cái xe đạp, buổi chiều bán nhiều hơn buổi sáng 5 cái. Buổi chiều bán được số xe đạp là
\choice
{$12$ cái}
{$14$ cái}
{$15$ cái}
{\True $13$ cái}
\loigiai{
$8 + 5 = 13$ cái.\\
\textbf{Đáp án của câu hỏi là D}
}
\end{ex}

\begin{ex} %[2D2V14]
Đàn vịt có 16 con, đàn gà có ít hơn đàn vịt 7 con. Đàn gà có số con là
\choice
{$8$ con}
{$10$ con}
{\True $9$ con}
{$11$ con}
\loigiai{
$16 - 7 = 9$ con.\\
\textbf{Đáp án của câu hỏi là C}
}
\end{ex}

\begin{ex} %[2D2V14]
Tìm $x$ biết: $x + 7 = 16$.
\choice
{$8$}
{$10$}
{$11$}
{\True $9$}
\loigiai{
$x = 16 - 7 = 9$.\\
\textbf{Đáp án của câu hỏi là D}
}
\end{ex}

\begin{ex} %[2D2V14]
Tìm $y$ biết: $15 - y = 6$.
\choice
{$8$}
{\True $9$}
{$10$}
{$7$}
\loigiai{
$y = 15 - 6 = 9$.\\
\textbf{Đáp án của câu hỏi là B}
}
\end{ex}

\begin{ex} %[2D2V14]
Tính nhanh: $9 + 6 + 1$.
\choice
{$15$}
{$17$}
{\True $16$}
{$18$}
\loigiai{
$(9 + 1) + 6 = 16$.\\
\textbf{Đáp án của câu hỏi là C}
}
\end{ex}

\begin{ex} %[2D2C14]
Tìm một số biết rằng lấy số đó cộng với 8 rồi trừ đi 5 thì được kết quả là 12.
\choice
{$8$}
{\True $9$}
{$10$}
{$7$}
\loigiai{
$12 + 5 - 8 = 9$.\\
\textbf{Đáp án của câu hỏi là B}
}
\end{ex}

\textbf{Phần 2. Trắc nghiệm đúng sai}
\setcounterpageref{ex}{0}

\begin{ex} %[2D2H14]
Xét tính đúng sai của các phép toán trong phạm vi 20:
\choiceTF
{\True $9 + 9 = 18$.}
{\True $18 - 9 = 9$.}
{$14 - 7 = 8$.}
{\True $14 - 7 = 7$.}
\loigiai{
\begin{itemize}
    \item $9 + 9 = 18$: Đúng.
    \item $18 - 9 = 9$: Đúng.
    \item $14 - 7 = 7 \neq 8$. Ý 3 Sai, Ý 4 Đúng.
\end{itemize}
\begin{itemize}
    \item a) \textbf{Đúng}.
    \item b) \textbf{Đúng}.
    \item c) \textbf{Sai}.
    \item d) \textbf{Đúng}.
\end{itemize}
}
\end{ex}

\begin{ex} %[2D2V14]
Hà có 14 quyển truyện, Nam có ít hơn Hà 5 quyển truyện. Xét tính đúng sai:
\choiceTF
{\True Nam có 9 quyển truyện ($14 - 5 = 9$).}
{\True Cả hai bạn có tất cả 23 quyển truyện ($14 + 9 = 23$).}
{Hà có ít truyện hơn Nam.}
{\True Hà có nhiều truyện hơn Nam (14 quyển $> 9$ quyển).}
\loigiai{
\begin{itemize}
    \item Nam có 9 quyển: Đúng.
    \item Cả hai có 23 quyển: Đúng.
    \item Hà nhiều hơn (không phải ít hơn). Ý 3 Sai, Ý 4 Đúng.
\end{itemize}
\begin{itemize}
    \item a) \textbf{Đúng}.
    \item b) \textbf{Đúng}.
    \item c) \textbf{Sai}.
    \item d) \textbf{Đúng}.
\end{itemize}
}
\end{ex}

\textbf{Phần 3. Trả lời ngắn}
\setcounterpageref{ex}{0}

\begin{ex} %[2D2H14]
Tính: $8 + 9$.
\shortans{17}
\loigiai{
\textbf{Đáp số:} $17$\\
$8 + 9 = 17$.
}
\end{ex}

\begin{ex} %[2D2H14]
Tính: $16 - 8$.
\shortans{8}
\loigiai{
\textbf{Đáp số:} $8$\\
$16 - 8 = 8$.
}
\end{ex}

\begin{ex} %[2D2V14]
Tính giá trị biểu thức: $17 - 9 + 4$.
\shortans{12}
\loigiai{
\textbf{Đáp số:} $12$\\
$8 + 4 = 12$.
}
\end{ex}

\begin{ex} %[2D2V14]
Tìm $x$ biết: $x - 7 = 8$.
\shortans{15}
\loigiai{
\textbf{Đáp số:} $15$\\
$x = 8 + 7 = 15$.
}
\end{ex}

\begin{ex} %[2D2V14]
Một vườn hoa có 9 bông hoa cúc, số hoa hồng nhiều hơn hoa cúc là 7 bông. Hỏi có bao nhiêu bông hoa hồng?
\shortans{16}
\loigiai{
\textbf{Đáp số:} $16$\\
$9 + 7 = 16$ bông hoa.
}
\end{ex}

\begin{ex} %[2D2C14]
Có hai rổ táo, rổ thứ nhất có 8 quả táo. Nếu thêm 6 quả vào rổ thứ nhất thì số táo ở rổ thứ nhất nhiều hơn rổ thứ hai là 2 quả. Hỏi rổ thứ hai có bao nhiêu quả táo?
\shortans{12}
\loigiai{
\textbf{Đáp số:} $12$\\
Rổ 1 lúc sau: $8 + 6 = 14$ quả.\\
Rổ 2 có: $14 - 2 = 12$ quả.
}
\end{ex}
